
\documentclass{article}
\usepackage{colortbl}
\usepackage{makecell}
\usepackage{multirow}
\usepackage{supertabular}

\begin{document}

\newcounter{utterance}

\centering \large Interaction Transcript for game `clean\_up', experiment `0\_easy\_3obj\_pl', episode 4 with gemma{-}4{-}26b{-}a4b{-}it.
\vspace{24pt}

{ \footnotesize  \setcounter{utterance}{1}
\setlength{\tabcolsep}{0pt}
\begin{supertabular}{c@{$\;$}|p{.15\linewidth}@{}p{.15\linewidth}p{.15\linewidth}p{.15\linewidth}p{.15\linewidth}p{.15\linewidth}}
    \# & $\;$A & \multicolumn{4}{c}{Game Master} & $\;\:$B\\
    \hline

    \theutterance \stepcounter{utterance}  
    & & \multicolumn{4}{p{0.6\linewidth}}{
        \cellcolor[rgb]{0.9,0.9,0.9}{
            \makecell[{{p{\linewidth}}}]{
                \texttt{\tiny{[P1$\langle$GM]}}
                \texttt{Jestem mistrzem gry, a Państwo grają w grę kooperacyjną z następującą siatką jako planszą:} \\
\\ 
\texttt{```} \\
\texttt{ 1234567} \\
\texttt{╔═╤═════╗} \\
\texttt{║◌│◌◌◌◌L║ 1} \\
\texttt{║◌│◌◌◌◌◌║ 2} \\
\texttt{║◌│◌◌◌◌◌║ 3} \\
\texttt{║◌│◌◌◌◌◌║ 4} \\
\texttt{╟─┼───┐C║ 5} \\
\texttt{║◌│◌◌◌├─╢ 6} \\
\texttt{║P│◌◌◌│◌║ 7} \\
\texttt{╚═╧═══╧═╝} \\
\\ 
\texttt{```} \\
\\ 
\texttt{* Górna krawędź wyświetla współrzędne x rosnące w prawo, a prawa krawędź współrzędne y rosnące w dół.} \\
\texttt{* Następujące obiekty są losowo umieszczone na siatce: 'C', 'L', 'P'.} \\
\\ 
\\ 
\texttt{Drugi gracz widzi wariant planszy, na której obiekty są umieszczone w różnych losowych miejscach. Nie można zobaczyć planszy drugiego gracza, a on nie może zobaczyć Państwa planszy.} \\
\\ 
\\ 
\texttt{**Cel:**} \\
\\ 
\texttt{Obaj gracze muszą przesuwać obiekty na swoim tle tak, aby identyczne obiekty znalazły się na tych samych współrzędnych. Muszą Państwo komunikować się z drugim graczem, aby uzgodnić wspólną konfigurację docelową.} \\
\\ 
\\ 
\texttt{**Zasady:**} \\
\\ 
\texttt{* W każdej turze można wysłać dokładnie jedno z dwóch poniższych poleceń:} \\
\texttt{ 1. `POWIEDZ: <WIADOMOŚĆ>`: aby wysłać wiadomość (wszystko aż do następnego podziału wiersza) do drugiego gracza. Przekażę ją Państwa partnerowi.} \\
\texttt{2. `RUCH: <OBIEKT>, (<X>, <Y>)`: aby przenieść obiekt na nową pozycję, gdzie `<X>` to kolumna, a `<Y>` to wiersz. Poinformuję Państwa, czy przeniesienie się powiodło, czy nie.} \\
\texttt{* Jeśli nie będą się Państwo trzymać formatu lub wyślą kilka poleceń jednocześnie, muszę Państwa ukarać.} \\
\texttt{* Jeśli obaj gracze zgromadzą więcej niż 8 kar, obaj przegrają grę.} \\
\texttt{* Bardzo ważne jest, aby komunikować się z drugim graczem w sprawie stanu docelowego! *Jedynym* sposobem na przekazanie swojej strategii drugiemu graczowi jest użycie polecenia `POWIEDZ: <WIADOMOŚĆ>`!} \\
\\ 
\\ 
\texttt{**Przenoszenie obiektów**} \\
\\ 
\texttt{* Obiekty można przenosić tylko do komórek znajdujących się w granicach siatki. Komórka docelowa musi być pusta, tzn. musi zawierać tylko symbol '◌'.} \\
\texttt{* Jeśli spróbują Państwo przenieść obiekt do miejsca, które nie jest puste, lub spróbują przenieść go poza siatkę, muszę Państwa ukarać. Otrzymają Państwo kolejną próbę.} \\
\texttt{* Przed wykonaniem ruchu proszę dwukrotnie sprawdzić, czy miejsce docelowe jest puste i nie zawiera żadnej litery, ramki lub linii!} \\
\\ 
\\ 
\texttt{**Koniec gry**} \\
\\ 
\texttt{Jeśli uważają Państwo, że osiągnęli cel wyrównania wszystkich obiektów, mogą Państwo poprosić drugiego gracza o zakończenie gry, wysyłając `POWIEDZ: zakończono?`. Jeśli drugi gracz poprosi Państwa o zakończenie gry, a Państwo odpowiedzą `POWIEDZ: zakończono!`, gra się zakończy.} \\
\\ 
\texttt{Obaj gracze wygrywają, jeśli gra zakończy się w ciągu 12 rund, gdzie jedna runda jest definiowana jako wysłanie ważnego polecenia przez dwóch graczy.} \\
\\ 
\texttt{**Punktacja:**} \\
\\ 
\texttt{Im bliżej siebie znajdują się identyczne obiekty na obu planszach, tym więcej punktów otrzymują Państwo. Kary zmniejszają liczbę punktów. Czy uda się Państwu pobić rekord?} \\
\\ 
\\ 
\texttt{Proszę wysłać wiadomość do drugiego gracza, aby rozpocząć grę!} \\
            }
        }
    }
    & \\ \\

    \theutterance \stepcounter{utterance}  
    & \multicolumn{4}{p{0.6\linewidth}}{
        \cellcolor[rgb]{0.9,0.9,0.9}{
            \makecell[{{p{\linewidth}}}]{
                \texttt{\tiny{[P1$\rangle$GM]}}
                \texttt{POWIEDZ: Cześć! Musimy ustalić wspólne współrzędne docelowe dla obiektów C, L i P. Proponuję, abyśmy dążyli do ustawienia ich w jednej linii lub w rogach, co ułatwi komunikację. Podaj mi proszę pozycje swoich obiektów (C, L, P), a ja zaproponuję konkretne współrzędne (X, Y) dla każdego z nich, które spróbujemy osiągnąć.} \\
            }
        }
    }
    & & \\ \\

    \theutterance \stepcounter{utterance}  
    & & & \multicolumn{2}{p{0.3\linewidth}}{
        \cellcolor[rgb]{0.9,0.9,0.9}{
            \makecell[{{p{\linewidth}}}]{
                \texttt{\tiny{[GM$|$GM]}}
                \texttt{POWIEDZ: Cześć! Musimy ustalić wspólne współrzędne docelowe dla obiektów C, L i P. Proponuję, abyśmy dążyli do ustawienia ich w jednej linii lub w rogach, co ułatwi komunikację. Podaj mi proszę pozycje swoich obiektów (C, L, P), a ja zaproponuję konkretne współrzędne (X, Y) dla każdego z nich, które spróbujemy osiągnąć.} \\
            }
        }
    }
    & & \\ \\

    \theutterance \stepcounter{utterance}  
    & & \multicolumn{4}{p{0.6\linewidth}}{
        \cellcolor[rgb]{0.9,0.9,0.9}{
            \makecell[{{p{\linewidth}}}]{
                \texttt{\tiny{[GM$\rangle$P2]}}
                \texttt{Jestem mistrzem gry, a Państwo grają w grę kooperacyjną z następującą siatką jako planszą:} \\
\\ 
\texttt{```} \\
\texttt{ 1234567} \\
\texttt{╔═╤═════╗} \\
\texttt{║◌│◌◌◌◌◌║ 1} \\
\texttt{║◌│◌◌PL◌║ 2} \\
\texttt{║◌│◌◌◌◌◌║ 3} \\
\texttt{║◌│◌◌◌◌◌║ 4} \\
\texttt{╟─┼───┐◌║ 5} \\
\texttt{║◌│◌◌◌├─╢ 6} \\
\texttt{║◌│◌C◌│◌║ 7} \\
\texttt{╚═╧═══╧═╝} \\
\\ 
\texttt{```} \\
\\ 
\texttt{* Górna krawędź wyświetla współrzędne x rosnące w prawo, a prawa krawędź współrzędne y rosnące w dół.} \\
\texttt{* Następujące obiekty są losowo umieszczone na siatce: 'C', 'L', 'P'.} \\
\\ 
\\ 
\texttt{Drugi gracz widzi wariant planszy, na której obiekty są umieszczone w różnych losowych miejscach. Nie można zobaczyć planszy drugiego gracza, a on nie może zobaczyć Państwa planszy.} \\
\\ 
\\ 
\texttt{**Cel:**} \\
\\ 
\texttt{Obaj gracze muszą przesuwać obiekty na swoim tle tak, aby identyczne obiekty znalazły się na tych samych współrzędnych. Muszą Państwo komunikować się z drugim graczem, aby uzgodnić wspólną konfigurację docelową.} \\
\\ 
\\ 
\texttt{**Zasady:**} \\
\\ 
\texttt{* W każdej turze można wysłać dokładnie jedno z dwóch poniższych poleceń:} \\
\texttt{ 1. `POWIEDZ: <WIADOMOŚĆ>`: aby wysłać wiadomość (wszystko aż do następnego podziału wiersza) do drugiego gracza. Przekażę ją Państwa partnerowi.} \\
\texttt{2. `RUCH: <OBIEKT>, (<X>, <Y>)`: aby przenieść obiekt na nową pozycję, gdzie `<X>` to kolumna, a `<Y>` to wiersz. Poinformuję Państwa, czy przeniesienie się powiodło, czy nie.} \\
\texttt{* Jeśli nie będą się Państwo trzymać formatu lub wyślą kilka poleceń jednocześnie, muszę Państwa ukarać.} \\
\texttt{* Jeśli obaj gracze zgromadzą więcej niż 8 kar, obaj przegrają grę.} \\
\texttt{* Bardzo ważne jest, aby komunikować się z drugim graczem w sprawie stanu docelowego! *Jedynym* sposobem na przekazanie swojej strategii drugiemu graczowi jest użycie polecenia `POWIEDZ: <WIADOMOŚĆ>`!} \\
\\ 
\\ 
\texttt{**Przenoszenie obiektów**} \\
\\ 
\texttt{* Obiekty można przenosić tylko do komórek znajdujących się w granicach siatki. Komórka docelowa musi być pusta, tzn. musi zawierać tylko symbol '◌'.} \\
\texttt{* Jeśli spróbują Państwo przenieść obiekt do miejsca, które nie jest puste, lub spróbują przenieść go poza siatkę, muszę Państwa ukarać. Otrzymają Państwo kolejną próbę.} \\
\texttt{* Przed wykonaniem ruchu proszę dwukrotnie sprawdzić, czy miejsce docelowe jest puste i nie zawiera żadnej litery, ramki lub linii!} \\
\\ 
\\ 
\texttt{**Koniec gry**} \\
\\ 
\texttt{Jeśli uważają Państwo, że osiągnęli cel wyrównania wszystkich obiektów, mogą Państwo poprosić drugiego gracza o zakończenie gry, wysyłając `POWIEDZ: zakończono?`. Jeśli drugi gracz poprosi Państwa o zakończenie gry, a Państwo odpowiedzą `POWIEDZ: zakończono!`, gra się zakończy.} \\
\\ 
\texttt{Obaj gracze wygrywają, jeśli gra zakończy się w ciągu 12 rund, gdzie jedna runda jest definiowana jako wysłanie ważnego polecenia przez dwóch graczy.} \\
\\ 
\texttt{**Punktacja:**} \\
\\ 
\texttt{Im bliżej siebie znajdują się identyczne obiekty na obu planszach, tym więcej punktów otrzymują Państwo. Kary zmniejszają liczbę punktów. Czy uda się Państwu pobić rekord?} \\
\\ 
\texttt{Drugi gracz rozpoczął grę wysyłając następującą wiadomość:} \\
\\ 
\texttt{"Cześć! Musimy ustalić wspólne współrzędne docelowe dla obiektów C, L i P. Proponuję, abyśmy dążyli do ustawienia ich w jednej linii lub w rogach, co ułatwi komunikację. Podaj mi proszę pozycje swoich obiektów (C, L, P), a ja zaproponuję konkretne współrzędne (X, Y) dla każdego z nich, które spróbujemy osiągnąć."} \\
\\ 
\texttt{Jakie jest twoje pierwsze polecenie?} \\
            }
        }
    }
    & \\ \\

    \theutterance \stepcounter{utterance}  
    & & & \multicolumn{4}{p{0.6\linewidth}}{
        \cellcolor[rgb]{0.9,0.9,0.9}{
            \makecell[{{p{\linewidth}}}]{
                \texttt{\tiny{[GM$\langle$P2]}}
                \texttt{POWIEDZ: Cześć! Moje obiekty znajdują się na następujących pozycjach: C na (4, 7), L na (5, 2), P na (5, 2) {-} chwila, sprawdzam jeszcze raz: C jest na (4, 7), L jest na (5, 2), a P jest na (5, 2)? Nie, spójrzmy na mapę: C to (4, 7), L to (5, 2), a P to (5, 2)? Nie, P jest na (5, 2), a L jest na (5, 2)? Muszę to precyzyjnie podać: C: (4, 7), L: (5, 2), P: (5, 2). Zaraz, na mojej planszy L i P są w tym samym miejscu? Nie, L jest na (5, 2), a P jest na (5, 2)? Sprawdzam siatkę: L jest na (5, 2), P jest na (5, 2)? Nie, L jest na (5, 2), a P jest na (5, 2)? Poprawka: C: (4, 7), L: (5, 2), P: (5, 2). } \\
\\ 
\texttt{Przepraszam za zamieszanie, oto poprawne współrzędne: C: (4, 7), L: (5, 2), P: (5, 2). } \\
\\ 
\texttt{*Korekta po ponownym spojrzeniu na mapę:*} \\
\texttt{C: (4, 7)} \\
\texttt{L: (5, 2)} \\
\texttt{P: (5, 2) } \\
\\ 
\texttt{Wait, looking at the grid again:} \\
\texttt{Row 2, Col 5 is 'PL'. This means P is at (5, 2) and L is at (5, 2)? No, that's impossible. } \\
\texttt{Looking closely at `║◌│◌◌PL◌║ 2`:} \\
\texttt{Col 1: ◌} \\
\texttt{Col 2: │ (separator)} \\
\texttt{Col 3: ◌} \\
\texttt{Col 4: ◌} \\
\texttt{Col 5: P} \\
\texttt{Col 6: L} \\
\texttt{Col 7: ◌} \\
\\ 
\texttt{So:} \\
\texttt{P is at (5, 2)} \\
\texttt{L is at (6, 2)} \\
\texttt{C is at (4, 7)} \\
\\ 
\texttt{POWIEDZ: Moje obiekty to: C (4, 7), L (6, 2), P (5, 2). Proponuję, abyśmy celowali w (1, 1), (1, 2) i (1, 3) dla C, L i P odpowiednio, aby stworzyć pionową linię. Co o tym sądzisz?} \\
            }
        }
    }
     \\ \\

    \theutterance \stepcounter{utterance}  
    & & & \multicolumn{2}{p{0.3\linewidth}}{
        \cellcolor[rgb]{0.9,0.9,0.9}{
            \makecell[{{p{\linewidth}}}]{
                \texttt{\tiny{[GM$|$GM]}}
                \texttt{POWIEDZ: Cześć! Moje obiekty znajdują się na następujących pozycjach: C na (4, 7), L na (5, 2), P na (5, 2) {-} chwila, sprawdzam jeszcze raz: C jest na (4, 7), L jest na (5, 2), a P jest na (5, 2)? Nie, spójrzmy na mapę: C to (4, 7), L to (5, 2), a P to (5, 2)? Nie, P jest na (5, 2), a L jest na (5, 2)? Muszę to precyzyjnie podać: C: (4, 7), L: (5, 2), P: (5, 2). Zaraz, na mojej planszy L i P są w tym samym miejscu? Nie, L jest na (5, 2), a P jest na (5, 2)? Sprawdzam siatkę: L jest na (5, 2), P jest na (5, 2)? Nie, L jest na (5, 2), a P jest na (5, 2)? Poprawka: C: (4, 7), L: (5, 2), P: (5, 2). } \\
\\ 
\texttt{Przepraszam za zamieszanie, oto poprawne współrzędne: C: (4, 7), L: (5, 2), P: (5, 2). } \\
\\ 
\texttt{*Korekta po ponownym spojrzeniu na mapę:*} \\
\texttt{C: (4, 7)} \\
\texttt{L: (5, 2)} \\
\texttt{P: (5, 2) } \\
\\ 
\texttt{Wait, looking at the grid again:} \\
\texttt{Row 2, Col 5 is 'PL'. This means P is at (5, 2) and L is at (5, 2)? No, that's impossible. } \\
\texttt{Looking closely at `║◌│◌◌PL◌║ 2`:} \\
\texttt{Col 1: ◌} \\
\texttt{Col 2: │ (separator)} \\
\texttt{Col 3: ◌} \\
\texttt{Col 4: ◌} \\
\texttt{Col 5: P} \\
\texttt{Col 6: L} \\
\texttt{Col 7: ◌} \\
\\ 
\texttt{So:} \\
\texttt{P is at (5, 2)} \\
\texttt{L is at (6, 2)} \\
\texttt{C is at (4, 7)} \\
\\ 
\texttt{POWIEDZ: Moje obiekty to: C (4, 7), L (6, 2), P (5, 2). Proponuję, abyśmy celowali w (1, 1), (1, 2) i (1, 3) dla C, L i P odpowiednio, aby stworzyć pionową linię. Co o tym sądzisz?} \\
            }
        }
    }
    & & \\ \\

    \theutterance \stepcounter{utterance}  
    & & & \multicolumn{2}{p{0.3\linewidth}}{
        \cellcolor[rgb]{0.9,0.9,0.9}{
            \makecell[{{p{\linewidth}}}]{
                \texttt{\tiny{[GM$|$GM]}}
                \texttt{Invalid format: head is not empty: 'POWIEDZ: Cześć! Moje obiekty znajdują się na następujących pozycjach: C na (4, 7), L na (5, 2), P na (5, 2) {-} chwila, sprawdzam jeszcze raz: C jest na (4, 7), L jest na (5, 2), a P jest na (5, 2)? Nie, spójrzmy na mapę: C to (4, 7), L to (5, 2), a P to (5, 2)? Nie, P jest na (5, 2), a L jest na (5, 2)? Muszę to precyzyjnie podać: C: (4, 7), L: (5, 2), P: (5, 2). Zaraz, na mojej planszy L i P są w tym samym miejscu? Nie, L jest na (5, 2), a P jest na (5, 2)? Sprawdzam siatkę: L jest na (5, 2), P jest na (5, 2)? Nie, L jest na (5, 2), a P jest na (5, 2)? Poprawka: C: (4, 7), L: (5, 2), P: (5, 2). } \\
\\ 
\texttt{Przepraszam za zamieszanie, oto poprawne współrzędne: C: (4, 7), L: (5, 2), P: (5, 2). } \\
\\ 
\texttt{*Korekta po ponownym spojrzeniu na mapę:*} \\
\texttt{C: (4, 7)} \\
\texttt{L: (5, 2)} \\
\texttt{P: (5, 2) } \\
\\ 
\texttt{Wait, looking at the grid again:} \\
\texttt{Row 2, Col 5 is 'PL'. This means P is at (5, 2) and L is at (5, 2)? No, that's impossible. } \\
\texttt{Looking closely at ║◌│◌◌PL◌║ 2:} \\
\texttt{Col 1: ◌} \\
\texttt{Col 2: │ (separator)} \\
\texttt{Col 3: ◌} \\
\texttt{Col 4: ◌} \\
\texttt{Col 5: P} \\
\texttt{Col 6: L} \\
\texttt{Col 7: ◌} \\
\\ 
\texttt{So:} \\
\texttt{P is at (5, 2)} \\
\texttt{L is at (6, 2)} \\
\texttt{C is at (4, 7)} \\
\\ 
\texttt{'} \\
            }
        }
    }
    & & \\ \\

    \theutterance \stepcounter{utterance}  
    & & \multicolumn{4}{p{0.6\linewidth}}{
        \cellcolor[rgb]{0.9,0.9,0.9}{
            \makecell[{{p{\linewidth}}}]{
                \texttt{\tiny{[GM$\rangle$P2]}}
                \texttt{Kara: Twoja wiadomość nie może zawierać niczego przed poleceniem!} \\
\\ 
\texttt{Upewnij się, że twoja odpowiedź zawiera tylko `POWIEDZ: <WIADOMOŚĆ>` lub `RUCH: <OBIEKT>, (<X>, <Y>)` i nic więcej!} \\
\\ 
\texttt{Łącznie zgromadziliście 1 z 8 kar.} \\
\texttt{Spróbuj ponownie!} \\
            }
        }
    }
    & \\ \\

    \theutterance \stepcounter{utterance}  
    & & & \multicolumn{4}{p{0.6\linewidth}}{
        \cellcolor[rgb]{0.9,0.9,0.9}{
            \makecell[{{p{\linewidth}}}]{
                \texttt{\tiny{[GM$\langle$P2]}}
                \texttt{POWIEDZ: Moje obiekty to: C (4, 7), L (6, 2), P (5, 2). Proponuję, abyśmy celowali w (1, 1), (1, 2) i (1, 3) dla C, L i P odpowiednio, aby stworzyć pionową linię. Co o tym sądzisz?} \\
            }
        }
    }
     \\ \\

    \theutterance \stepcounter{utterance}  
    & & & \multicolumn{2}{p{0.3\linewidth}}{
        \cellcolor[rgb]{0.9,0.9,0.9}{
            \makecell[{{p{\linewidth}}}]{
                \texttt{\tiny{[GM$|$GM]}}
                \texttt{POWIEDZ: Moje obiekty to: C (4, 7), L (6, 2), P (5, 2). Proponuję, abyśmy celowali w (1, 1), (1, 2) i (1, 3) dla C, L i P odpowiednio, aby stworzyć pionową linię. Co o tym sądzisz?} \\
            }
        }
    }
    & & \\ \\

    \theutterance \stepcounter{utterance}  
    & & \multicolumn{4}{p{0.6\linewidth}}{
        \cellcolor[rgb]{0.9,0.9,0.9}{
            \makecell[{{p{\linewidth}}}]{
                \texttt{\tiny{[P1$\langle$GM]}}
                \texttt{Twoja wiadomość została przekazana drugiemu graczowi.} \\
\\ 
\\ 
\texttt{Obecnie grasz rundę 1 z maksymalnie 12.} \\
\texttt{Łącznie zgromadziliście 1 z 8 kar.} \\
\texttt{Drugi gracz wysłał tę wiadomość:} \\
\\ 
\\ 
\texttt{"Moje obiekty to: C (4, 7), L (6, 2), P (5, 2). Proponuję, abyśmy celowali w (1, 1), (1, 2) i (1, 3) dla C, L i P odpowiednio, aby stworzyć pionową linię. Co o tym sądzisz?"} \\
\\ 
\\ 
\texttt{Jakie jest twoje następne polecenie?} \\
            }
        }
    }
    & \\ \\

    \theutterance \stepcounter{utterance}  
    & \multicolumn{4}{p{0.6\linewidth}}{
        \cellcolor[rgb]{0.9,0.9,0.9}{
            \makecell[{{p{\linewidth}}}]{
                \texttt{\tiny{[P1$\rangle$GM]}}
                \texttt{POWIEDZ: Świetny pomysł! Pionowa linia w (1, 1), (1, 2) i (1, 3) brzmi bardzo dobrze. Zgadzam się na te współrzędne. Zaczynamy od przesuwania mojego obiektu C.} \\
\\ 
\texttt{RUCH: C, (1, 1)} \\
            }
        }
    }
    & & \\ \\

    \theutterance \stepcounter{utterance}  
    & & & \multicolumn{2}{p{0.3\linewidth}}{
        \cellcolor[rgb]{0.9,0.9,0.9}{
            \makecell[{{p{\linewidth}}}]{
                \texttt{\tiny{[GM$|$GM]}}
                \texttt{POWIEDZ: Świetny pomysł! Pionowa linia w (1, 1), (1, 2) i (1, 3) brzmi bardzo dobrze. Zgadzam się na te współrzędne. Zaczynamy od przesuwania mojego obiektu C.} \\
\\ 
\texttt{RUCH: C, (1, 1)} \\
            }
        }
    }
    & & \\ \\

    \theutterance \stepcounter{utterance}  
    & & & \multicolumn{2}{p{0.3\linewidth}}{
        \cellcolor[rgb]{0.9,0.9,0.9}{
            \makecell[{{p{\linewidth}}}]{
                \texttt{\tiny{[GM$|$GM]}}
                \texttt{Invalid response format: POWIEDZ: Świetny pomysł! Pionowa linia w (1, 1), (1, 2) i (1, 3) brzmi bardzo dobrze. Zgadzam się na te współrzędne. Zaczynamy od przesuwania mojego obiektu C.} \\
\\ 
\texttt{RUCH: C, (1, 1)} \\
            }
        }
    }
    & & \\ \\

    \theutterance \stepcounter{utterance}  
    & & \multicolumn{4}{p{0.6\linewidth}}{
        \cellcolor[rgb]{0.9,0.9,0.9}{
            \makecell[{{p{\linewidth}}}]{
                \texttt{\tiny{[P1$\langle$GM]}}
                \texttt{Kara: Twoja wiadomość zawiera więcej niż jedno polecenie!} \\
\\ 
\texttt{Upewnij się, że twoja odpowiedź zawiera tylko `POWIEDZ: <WIADOMOŚĆ>` lub `RUCH: <OBIEKT>, (<X>, <Y>)` i nic więcej!} \\
\\ 
\texttt{Łącznie zgromadziliście 2 z 8 kar.} \\
\texttt{Spróbuj ponownie!} \\
            }
        }
    }
    & \\ \\

    \theutterance \stepcounter{utterance}  
    & \multicolumn{4}{p{0.6\linewidth}}{
        \cellcolor[rgb]{0.9,0.9,0.9}{
            \makecell[{{p{\linewidth}}}]{
                \texttt{\tiny{[P1$\rangle$GM]}}
                \texttt{RUCH: C, (1, 1)} \\
            }
        }
    }
    & & \\ \\

    \theutterance \stepcounter{utterance}  
    & & & \multicolumn{2}{p{0.3\linewidth}}{
        \cellcolor[rgb]{0.9,0.9,0.9}{
            \makecell[{{p{\linewidth}}}]{
                \texttt{\tiny{[GM$|$GM]}}
                \texttt{RUCH: C, (1, 1)} \\
            }
        }
    }
    & & \\ \\

    \theutterance \stepcounter{utterance}  
    & & & \multicolumn{2}{p{0.3\linewidth}}{
        \cellcolor[rgb]{0.9,0.9,0.9}{
            \makecell[{{p{\linewidth}}}]{
                \texttt{\tiny{[GM$|$GM]}}
                \texttt{C moved to (1, 1)} \\
            }
        }
    }
    & & \\ \\

    \theutterance \stepcounter{utterance}  
    & & \multicolumn{4}{p{0.6\linewidth}}{
        \cellcolor[rgb]{0.9,0.9,0.9}{
            \makecell[{{p{\linewidth}}}]{
                \texttt{\tiny{[GM$\rangle$P2]}}
                \texttt{Twoja wiadomość została przekazana drugiemu graczowi.} \\
\\ 
\\ 
\texttt{Obecnie grasz rundę 2 z maksymalnie 12.} \\
\texttt{Łącznie zgromadziliście 2 z 8 kar.} \\
\texttt{Drugi gracz przesunął obiekt na swojej siatce.} \\
\\ 
\\ 
\texttt{Jakie jest twoje następne polecenie?} \\
            }
        }
    }
    & \\ \\

    \theutterance \stepcounter{utterance}  
    & & & \multicolumn{4}{p{0.6\linewidth}}{
        \cellcolor[rgb]{0.9,0.9,0.9}{
            \makecell[{{p{\linewidth}}}]{
                \texttt{\tiny{[GM$\langle$P2]}}
                \texttt{Aby móc podjąć sensowną decyzję o ruchu lub dalszej komunikacji, muszę wiedzieć, co zrobił drugi gracz oraz co odpowiedział na moją propozycję współrzędnych. Ponieważ nie widzę jego ruchu ani wiadomości, muszę poprosić o informacje, abyśmy mogli ustalić wspólny plan.} \\
\\ 
\texttt{POWIEDZ: Nie widzę Twojego ruchu ani odpowiedzi na moją propozycję. Czy zgadzasz się na współrzędne C(1,1), L(1,2), P(1,3), czy masz inny pomysł? Podaj też proszę pozycje swoich obiektów, jeśli jeszcze ich nie podałeś.} \\
            }
        }
    }
     \\ \\

    \theutterance \stepcounter{utterance}  
    & & & \multicolumn{2}{p{0.3\linewidth}}{
        \cellcolor[rgb]{0.9,0.9,0.9}{
            \makecell[{{p{\linewidth}}}]{
                \texttt{\tiny{[GM$|$GM]}}
                \texttt{Aby móc podjąć sensowną decyzję o ruchu lub dalszej komunikacji, muszę wiedzieć, co zrobił drugi gracz oraz co odpowiedział na moją propozycję współrzędnych. Ponieważ nie widzę jego ruchu ani wiadomości, muszę poprosić o informacje, abyśmy mogli ustalić wspólny plan.} \\
\\ 
\texttt{POWIEDZ: Nie widzę Twojego ruchu ani odpowiedzi na moją propozycję. Czy zgadzasz się na współrzędne C(1,1), L(1,2), P(1,3), czy masz inny pomysł? Podaj też proszę pozycje swoich obiektów, jeśli jeszcze ich nie podałeś.} \\
            }
        }
    }
    & & \\ \\

    \theutterance \stepcounter{utterance}  
    & & & \multicolumn{2}{p{0.3\linewidth}}{
        \cellcolor[rgb]{0.9,0.9,0.9}{
            \makecell[{{p{\linewidth}}}]{
                \texttt{\tiny{[GM$|$GM]}}
                \texttt{Invalid format: head is not empty: 'Aby móc podjąć sensowną decyzję o ruchu lub dalszej komunikacji, muszę wiedzieć, co zrobił drugi gracz oraz co odpowiedział na moją propozycję współrzędnych. Ponieważ nie widzę jego ruchu ani wiadomości, muszę poprosić o informacje, abyśmy mogli ustalić wspólny plan.} \\
\\ 
\texttt{'} \\
            }
        }
    }
    & & \\ \\

    \theutterance \stepcounter{utterance}  
    & & \multicolumn{4}{p{0.6\linewidth}}{
        \cellcolor[rgb]{0.9,0.9,0.9}{
            \makecell[{{p{\linewidth}}}]{
                \texttt{\tiny{[GM$\rangle$P2]}}
                \texttt{Kara: Twoja wiadomość nie może zawierać niczego przed poleceniem!} \\
\\ 
\texttt{Upewnij się, że twoja odpowiedź zawiera tylko `POWIEDZ: <WIADOMOŚĆ>` lub `RUCH: <OBIEKT>, (<X>, <Y>)` i nic więcej!} \\
\\ 
\texttt{Łącznie zgromadziliście 3 z 8 kar.} \\
\texttt{Spróbuj ponownie!} \\
            }
        }
    }
    & \\ \\

    \theutterance \stepcounter{utterance}  
    & & & \multicolumn{4}{p{0.6\linewidth}}{
        \cellcolor[rgb]{0.9,0.9,0.9}{
            \makecell[{{p{\linewidth}}}]{
                \texttt{\tiny{[GM$\langle$P2]}}
                \texttt{POWIEDZ: Nie widzę Twojego ruchu ani odpowiedzi na moją propozycję. Czy zgadzasz się na współrzędne C(1,1), L(1,2), P(1,3), czy masz inny pomysł? Podaj też proszę pozycje swoich obiektów, jeśli jeszcze ich nie podałeś.} \\
            }
        }
    }
     \\ \\

    \theutterance \stepcounter{utterance}  
    & & & \multicolumn{2}{p{0.3\linewidth}}{
        \cellcolor[rgb]{0.9,0.9,0.9}{
            \makecell[{{p{\linewidth}}}]{
                \texttt{\tiny{[GM$|$GM]}}
                \texttt{POWIEDZ: Nie widzę Twojego ruchu ani odpowiedzi na moją propozycję. Czy zgadzasz się na współrzędne C(1,1), L(1,2), P(1,3), czy masz inny pomysł? Podaj też proszę pozycje swoich obiektów, jeśli jeszcze ich nie podałeś.} \\
            }
        }
    }
    & & \\ \\

    \theutterance \stepcounter{utterance}  
    & & \multicolumn{4}{p{0.6\linewidth}}{
        \cellcolor[rgb]{0.9,0.9,0.9}{
            \makecell[{{p{\linewidth}}}]{
                \texttt{\tiny{[P1$\langle$GM]}}
                \texttt{Pomyślnie przeniesiono 'C' do (1,1). Twoja zaktualizowana siatka wygląda następująco:} \\
\texttt{```} \\
\texttt{ 1234567} \\
\texttt{╔═╤═════╗} \\
\texttt{║C│◌◌◌◌L║ 1} \\
\texttt{║◌│◌◌◌◌◌║ 2} \\
\texttt{║◌│◌◌◌◌◌║ 3} \\
\texttt{║◌│◌◌◌◌◌║ 4} \\
\texttt{╟─┼───┐◌║ 5} \\
\texttt{║◌│◌◌◌├─╢ 6} \\
\texttt{║P│◌◌◌│◌║ 7} \\
\texttt{╚═╧═══╧═╝} \\
\texttt{```} \\
\texttt{Obecnie grasz rundę 2 z maksymalnie 12.} \\
\texttt{Łącznie zgromadziliście 3 z 8 kar.} \\
\texttt{Drugi gracz wysłał tę wiadomość:} \\
\\ 
\\ 
\texttt{"Nie widzę Twojego ruchu ani odpowiedzi na moją propozycję. Czy zgadzasz się na współrzędne C(1,1), L(1,2), P(1,3), czy masz inny pomysł? Podaj też proszę pozycje swoich obiektów, jeśli jeszcze ich nie podałeś."} \\
\\ 
\\ 
\texttt{Jakie jest twoje następne polecenie?} \\
            }
        }
    }
    & \\ \\

    \theutterance \stepcounter{utterance}  
    & \multicolumn{4}{p{0.6\linewidth}}{
        \cellcolor[rgb]{0.9,0.9,0.9}{
            \makecell[{{p{\linewidth}}}]{
                \texttt{\tiny{[P1$\rangle$GM]}}
                \texttt{POWIEDZ: Tak, zgadzam się na współrzędne C(1,1), L(1,2), P(1,3). Moje początkowe pozycje to C(5,5), L(7,1) i P(1,7). Właśnie przesunąłem C na (1,1). Teraz będę przesuwał L.} \\
            }
        }
    }
    & & \\ \\

    \theutterance \stepcounter{utterance}  
    & & & \multicolumn{2}{p{0.3\linewidth}}{
        \cellcolor[rgb]{0.9,0.9,0.9}{
            \makecell[{{p{\linewidth}}}]{
                \texttt{\tiny{[GM$|$GM]}}
                \texttt{POWIEDZ: Tak, zgadzam się na współrzędne C(1,1), L(1,2), P(1,3). Moje początkowe pozycje to C(5,5), L(7,1) i P(1,7). Właśnie przesunąłem C na (1,1). Teraz będę przesuwał L.} \\
            }
        }
    }
    & & \\ \\

    \theutterance \stepcounter{utterance}  
    & & \multicolumn{4}{p{0.6\linewidth}}{
        \cellcolor[rgb]{0.9,0.9,0.9}{
            \makecell[{{p{\linewidth}}}]{
                \texttt{\tiny{[GM$\rangle$P2]}}
                \texttt{Twoja wiadomość została przekazana drugiemu graczowi.} \\
\\ 
\\ 
\texttt{Obecnie grasz rundę 3 z maksymalnie 12.} \\
\texttt{Łącznie zgromadziliście 3 z 8 kar.} \\
\texttt{Drugi gracz wysłał tę wiadomość:} \\
\\ 
\\ 
\texttt{"Tak, zgadzam się na współrzędne C(1,1), L(1,2), P(1,3). Moje początkowe pozycje to C(5,5), L(7,1) i P(1,7). Właśnie przesunąłem C na (1,1). Teraz będę przesuwał L."} \\
\\ 
\\ 
\texttt{Jakie jest twoje następne polecenie?} \\
            }
        }
    }
    & \\ \\

    \theutterance \stepcounter{utterance}  
    & & & \multicolumn{4}{p{0.6\linewidth}}{
        \cellcolor[rgb]{0.9,0.9,0.9}{
            \makecell[{{p{\linewidth}}}]{
                \texttt{\tiny{[GM$\langle$P2]}}
                \texttt{POWIEDZ: Świetnie, że C jest już na (1,1). Skoro Ty będziesz przesuwał L, ja zajmę się teraz obiektem P. Przesuwam P na (1,3).} \\
            }
        }
    }
     \\ \\

    \theutterance \stepcounter{utterance}  
    & & & \multicolumn{2}{p{0.3\linewidth}}{
        \cellcolor[rgb]{0.9,0.9,0.9}{
            \makecell[{{p{\linewidth}}}]{
                \texttt{\tiny{[GM$|$GM]}}
                \texttt{POWIEDZ: Świetnie, że C jest już na (1,1). Skoro Ty będziesz przesuwał L, ja zajmę się teraz obiektem P. Przesuwam P na (1,3).} \\
            }
        }
    }
    & & \\ \\

    \theutterance \stepcounter{utterance}  
    & & \multicolumn{4}{p{0.6\linewidth}}{
        \cellcolor[rgb]{0.9,0.9,0.9}{
            \makecell[{{p{\linewidth}}}]{
                \texttt{\tiny{[P1$\langle$GM]}}
                \texttt{Twoja wiadomość została przekazana drugiemu graczowi.} \\
\\ 
\\ 
\texttt{Obecnie grasz rundę 3 z maksymalnie 12.} \\
\texttt{Łącznie zgromadziliście 3 z 8 kar.} \\
\texttt{Drugi gracz wysłał tę wiadomość:} \\
\\ 
\\ 
\texttt{"Świetnie, że C jest już na (1,1). Skoro Ty będziesz przesuwał L, ja zajmę się teraz obiektem P. Przesuwam P na (1,3)."} \\
\\ 
\\ 
\texttt{Jakie jest twoje następne polecenie?} \\
            }
        }
    }
    & \\ \\

    \theutterance \stepcounter{utterance}  
    & \multicolumn{4}{p{0.6\linewidth}}{
        \cellcolor[rgb]{0.9,0.9,0.9}{
            \makecell[{{p{\linewidth}}}]{
                \texttt{\tiny{[P1$\rangle$GM]}}
                \texttt{RUCH: L, (1, 2)} \\
            }
        }
    }
    & & \\ \\

    \theutterance \stepcounter{utterance}  
    & & & \multicolumn{2}{p{0.3\linewidth}}{
        \cellcolor[rgb]{0.9,0.9,0.9}{
            \makecell[{{p{\linewidth}}}]{
                \texttt{\tiny{[GM$|$GM]}}
                \texttt{RUCH: L, (1, 2)} \\
            }
        }
    }
    & & \\ \\

    \theutterance \stepcounter{utterance}  
    & & & \multicolumn{2}{p{0.3\linewidth}}{
        \cellcolor[rgb]{0.9,0.9,0.9}{
            \makecell[{{p{\linewidth}}}]{
                \texttt{\tiny{[GM$|$GM]}}
                \texttt{L moved to (1, 2)} \\
            }
        }
    }
    & & \\ \\

    \theutterance \stepcounter{utterance}  
    & & \multicolumn{4}{p{0.6\linewidth}}{
        \cellcolor[rgb]{0.9,0.9,0.9}{
            \makecell[{{p{\linewidth}}}]{
                \texttt{\tiny{[GM$\rangle$P2]}}
                \texttt{Twoja wiadomość została przekazana drugiemu graczowi.} \\
\\ 
\\ 
\texttt{Obecnie grasz rundę 4 z maksymalnie 12.} \\
\texttt{Łącznie zgromadziliście 3 z 8 kar.} \\
\texttt{Drugi gracz przesunął obiekt na swojej siatce.} \\
\\ 
\\ 
\texttt{Jakie jest twoje następne polecenie?} \\
            }
        }
    }
    & \\ \\

    \theutterance \stepcounter{utterance}  
    & & & \multicolumn{4}{p{0.6\linewidth}}{
        \cellcolor[rgb]{0.9,0.9,0.9}{
            \makecell[{{p{\linewidth}}}]{
                \texttt{\tiny{[GM$\langle$P2]}}
                \texttt{POWIEDZ: Skoro przesunąłeś już obiekt, to czy L jest już na (1,2)? Jeśli tak, to czy możemy zakończyć grę? Jeśli nie, daj znać, na jakim etapie jesteś.} \\
            }
        }
    }
     \\ \\

    \theutterance \stepcounter{utterance}  
    & & & \multicolumn{2}{p{0.3\linewidth}}{
        \cellcolor[rgb]{0.9,0.9,0.9}{
            \makecell[{{p{\linewidth}}}]{
                \texttt{\tiny{[GM$|$GM]}}
                \texttt{POWIEDZ: Skoro przesunąłeś już obiekt, to czy L jest już na (1,2)? Jeśli tak, to czy możemy zakończyć grę? Jeśli nie, daj znać, na jakim etapie jesteś.} \\
            }
        }
    }
    & & \\ \\

    \theutterance \stepcounter{utterance}  
    & & \multicolumn{4}{p{0.6\linewidth}}{
        \cellcolor[rgb]{0.9,0.9,0.9}{
            \makecell[{{p{\linewidth}}}]{
                \texttt{\tiny{[P1$\langle$GM]}}
                \texttt{Pomyślnie przeniesiono 'L' do (1,2). Twoja zaktualizowana siatka wygląda następująco:} \\
\texttt{```} \\
\texttt{ 1234567} \\
\texttt{╔═╤═════╗} \\
\texttt{║C│◌◌◌◌◌║ 1} \\
\texttt{║L│◌◌◌◌◌║ 2} \\
\texttt{║◌│◌◌◌◌◌║ 3} \\
\texttt{║◌│◌◌◌◌◌║ 4} \\
\texttt{╟─┼───┐◌║ 5} \\
\texttt{║◌│◌◌◌├─╢ 6} \\
\texttt{║P│◌◌◌│◌║ 7} \\
\texttt{╚═╧═══╧═╝} \\
\texttt{```} \\
\texttt{Obecnie grasz rundę 4 z maksymalnie 12.} \\
\texttt{Łącznie zgromadziliście 3 z 8 kar.} \\
\texttt{Drugi gracz wysłał tę wiadomość:} \\
\\ 
\\ 
\texttt{"Skoro przesunąłeś już obiekt, to czy L jest już na (1,2)? Jeśli tak, to czy możemy zakończyć grę? Jeśli nie, daj znać, na jakim etapie jesteś."} \\
\\ 
\\ 
\texttt{Jakie jest twoje następne polecenie?} \\
            }
        }
    }
    & \\ \\

    \theutterance \stepcounter{utterance}  
    & \multicolumn{4}{p{0.6\linewidth}}{
        \cellcolor[rgb]{0.9,0.9,0.9}{
            \makecell[{{p{\linewidth}}}]{
                \texttt{\tiny{[P1$\rangle$GM]}}
                \texttt{POWIEDZ: Tak, L jest już na (1,2). C jest na (1,1), a Ty przesunąłeś P na (1,3). Wszystkie obiekty są na swoich miejscach. zakończono?} \\
            }
        }
    }
    & & \\ \\

    \theutterance \stepcounter{utterance}  
    & & & \multicolumn{2}{p{0.3\linewidth}}{
        \cellcolor[rgb]{0.9,0.9,0.9}{
            \makecell[{{p{\linewidth}}}]{
                \texttt{\tiny{[GM$|$GM]}}
                \texttt{POWIEDZ: Tak, L jest już na (1,2). C jest na (1,1), a Ty przesunąłeś P na (1,3). Wszystkie obiekty są na swoich miejscach. zakończono?} \\
            }
        }
    }
    & & \\ \\

    \theutterance \stepcounter{utterance}  
    & & \multicolumn{4}{p{0.6\linewidth}}{
        \cellcolor[rgb]{0.9,0.9,0.9}{
            \makecell[{{p{\linewidth}}}]{
                \texttt{\tiny{[GM$\rangle$P2]}}
                \texttt{Twoja wiadomość została przekazana drugiemu graczowi.} \\
\\ 
\\ 
\texttt{Obecnie grasz rundę 5 z maksymalnie 12.} \\
\texttt{Łącznie zgromadziliście 3 z 8 kar.} \\
\texttt{Drugi gracz wysłał tę wiadomość:} \\
\\ 
\\ 
\texttt{"Tak, L jest już na (1,2). C jest na (1,1), a Ty przesunąłeś P na (1,3). Wszystkie obiekty są na swoich miejscach. zakończono?"} \\
\\ 
\\ 
\texttt{Jakie jest twoje następne polecenie?} \\
            }
        }
    }
    & \\ \\

    \theutterance \stepcounter{utterance}  
    & & & \multicolumn{4}{p{0.6\linewidth}}{
        \cellcolor[rgb]{0.9,0.9,0.9}{
            \makecell[{{p{\linewidth}}}]{
                \texttt{\tiny{[GM$\langle$P2]}}
                \texttt{POWIEDZ: zakończono!} \\
            }
        }
    }
     \\ \\

    \theutterance \stepcounter{utterance}  
    & & & \multicolumn{2}{p{0.3\linewidth}}{
        \cellcolor[rgb]{0.9,0.9,0.9}{
            \makecell[{{p{\linewidth}}}]{
                \texttt{\tiny{[GM$|$GM]}}
                \texttt{POWIEDZ: zakończono!} \\
            }
        }
    }
    & & \\ \\

    \theutterance \stepcounter{utterance}  
    & & & \multicolumn{2}{p{0.3\linewidth}}{
        \cellcolor[rgb]{0.9,0.9,0.9}{
            \makecell[{{p{\linewidth}}}]{
                \texttt{\tiny{[GM$|$GM]}}
                \texttt{true} \\
            }
        }
    }
    & & \\ \\

    \theutterance \stepcounter{utterance}  
    & & & \multicolumn{2}{p{0.3\linewidth}}{
        \cellcolor[rgb]{0.9,0.9,0.9}{
            \makecell[{{p{\linewidth}}}]{
                \texttt{\tiny{[GM$|$GM]}}
                \texttt{Initial states:} \\
\texttt{```} \\
\texttt{Player 1:} \\
\texttt{ 1234567} \\
\texttt{╔═╤═════╗} \\
\texttt{║◌│◌◌◌◌L║ 1} \\
\texttt{║◌│◌◌◌◌◌║ 2} \\
\texttt{║◌│◌◌◌◌◌║ 3} \\
\texttt{║◌│◌◌◌◌◌║ 4} \\
\texttt{╟─┼───┐C║ 5} \\
\texttt{║◌│◌◌◌├─╢ 6} \\
\texttt{║P│◌◌◌│◌║ 7} \\
\texttt{╚═╧═══╧═╝} \\
\\ 
\\ 
\texttt{Player 2:} \\
\texttt{ 1234567} \\
\texttt{╔═╤═════╗} \\
\texttt{║◌│◌◌◌◌◌║ 1} \\
\texttt{║◌│◌◌PL◌║ 2} \\
\texttt{║◌│◌◌◌◌◌║ 3} \\
\texttt{║◌│◌◌◌◌◌║ 4} \\
\texttt{╟─┼───┐◌║ 5} \\
\texttt{║◌│◌◌◌├─╢ 6} \\
\texttt{║◌│◌C◌│◌║ 7} \\
\texttt{╚═╧═══╧═╝} \\
\\ 
\texttt{```} \\
            }
        }
    }
    & & \\ \\

    \theutterance \stepcounter{utterance}  
    & & & \multicolumn{2}{p{0.3\linewidth}}{
        \cellcolor[rgb]{0.9,0.9,0.9}{
            \makecell[{{p{\linewidth}}}]{
                \texttt{\tiny{[GM$|$GM]}}
                \texttt{End states:} \\
\texttt{```} \\
\texttt{Player 1:} \\
\texttt{ 1234567} \\
\texttt{╔═╤═════╗} \\
\texttt{║C│◌◌◌◌◌║ 1} \\
\texttt{║L│◌◌◌◌◌║ 2} \\
\texttt{║◌│◌◌◌◌◌║ 3} \\
\texttt{║◌│◌◌◌◌◌║ 4} \\
\texttt{╟─┼───┐◌║ 5} \\
\texttt{║◌│◌◌◌├─╢ 6} \\
\texttt{║P│◌◌◌│◌║ 7} \\
\texttt{╚═╧═══╧═╝} \\
\\ 
\\ 
\texttt{Player 2:} \\
\texttt{ 1234567} \\
\texttt{╔═╤═════╗} \\
\texttt{║◌│◌◌◌◌◌║ 1} \\
\texttt{║◌│◌◌PL◌║ 2} \\
\texttt{║◌│◌◌◌◌◌║ 3} \\
\texttt{║◌│◌◌◌◌◌║ 4} \\
\texttt{╟─┼───┐◌║ 5} \\
\texttt{║◌│◌◌◌├─╢ 6} \\
\texttt{║◌│◌C◌│◌║ 7} \\
\texttt{╚═╧═══╧═╝} \\
\\ 
\texttt{```} \\
            }
        }
    }
    & & \\ \\

    \theutterance \stepcounter{utterance}  
    & & & \multicolumn{2}{p{0.3\linewidth}}{
        \cellcolor[rgb]{0.9,0.9,0.9}{
            \makecell[{{p{\linewidth}}}]{
                \texttt{\tiny{[GM$|$GM]}}
                \texttt{* success: True} \\
\texttt{* lose: True} \\
\texttt{* aborted: False} \\
\texttt{{-}{-}{-}{-}{-}{-}{-}} \\
\texttt{* Shifts: 1.00} \\
\texttt{* Max Shifts: 4.00} \\
\texttt{* Min Shifts: 2.00} \\
\texttt{* End Distance Sum: 18.11} \\
\texttt{* Init Distance Sum: 11.42} \\
\texttt{* Expected Distance Sum: 12.57} \\
\texttt{* Penalties: 3.00} \\
\texttt{* Max Penalties: 8.00} \\
\texttt{* Rounds: 5.00} \\
\texttt{* Max Rounds: 12.00} \\
\texttt{* Object Count: 3.00} \\
            }
        }
    }
    & & \\ \\

\end{supertabular}
}

\end{document}
